% Options for packages loaded elsewhere
\PassOptionsToPackage{unicode}{hyperref}
\PassOptionsToPackage{hyphens}{url}
%
\documentclass[
  10pt,
]{article}
\usepackage{amsmath,amssymb}
\usepackage{iftex}
\ifPDFTeX
  \usepackage[T1]{fontenc}
  \usepackage[utf8]{inputenc}
  \usepackage{textcomp} % provide euro and other symbols
\else % if luatex or xetex
  \usepackage{unicode-math} % this also loads fontspec
  \defaultfontfeatures{Scale=MatchLowercase}
  \defaultfontfeatures[\rmfamily]{Ligatures=TeX,Scale=1}
\fi
\usepackage{lmodern}
\ifPDFTeX\else
  % xetex/luatex font selection
\fi
% Use upquote if available, for straight quotes in verbatim environments
\IfFileExists{upquote.sty}{\usepackage{upquote}}{}
\IfFileExists{microtype.sty}{% use microtype if available
  \usepackage[]{microtype}
  \UseMicrotypeSet[protrusion]{basicmath} % disable protrusion for tt fonts
}{}
\makeatletter
\@ifundefined{KOMAClassName}{% if non-KOMA class
  \IfFileExists{parskip.sty}{%
    \usepackage{parskip}
  }{% else
    \setlength{\parindent}{0pt}
    \setlength{\parskip}{6pt plus 2pt minus 1pt}}
}{% if KOMA class
  \KOMAoptions{parskip=half}}
\makeatother
\usepackage{xcolor}
\usepackage[margin=24mm]{geometry}
\setlength{\emergencystretch}{3em} % prevent overfull lines
\providecommand{\tightlist}{%
  \setlength{\itemsep}{0pt}\setlength{\parskip}{0pt}}
\setcounter{secnumdepth}{-\maxdimen} % remove section numbering
\ifLuaTeX
  \usepackage{selnolig}  % disable illegal ligatures
\fi
\usepackage{bookmark}
\IfFileExists{xurl.sty}{\usepackage{xurl}}{} % add URL line breaks if available
\urlstyle{same}
\hypersetup{
  hidelinks,
  pdfcreator={LaTeX via pandoc}}

\author{}
\date{}

\begin{document}

\section{Exact edge general position in cyclic run-avoiding cubes for
every k \textgreater=
3}\label{exact-edge-general-position-in-cyclic-run-avoiding-cubes-for-every-k-3}

Ryutaro Yonezu (Independent Researcher)\\
Research No.78 · Preprint version 1.0 · 20 November 2026

\textbf{Status.} AI-assisted, not human peer reviewed. THEOREM /
UNIQUENESS: CONDITIONAL PASS in the independent AI audit. NOVELTY
INCONCLUSIVE; REFERENCES PARTIAL. Copyright © 2026 Ryutaro Yonezu. All
rights reserved.

\subsection{Editorial scope of this combined
edition}\label{editorial-scope-of-this-combined-edition}

This edition contains the complete, unchanged mathematical prose of
frozen v0.2 and its parent v0.1. The v0.2 supplement establishes the
stated result for every integer n \textgreater= k \textgreater= 3; its
proofs for k \textgreater= 4 depend on the full parent v0.1 §§1--5,
reproduced below. The parent's unresolved k=3 statement is a historical
record of its earlier stage and is resolved by the later supplement. The
k \textgreater= 4 proof is analytic. The k = 3 proof uses exact integer
computation for 1,296 required transfer-matrix bases and an exhaustive
n=4 equality check, together with an analytic strong induction. This
edition makes no assertion for k=2 or n\textless k.

The independent auditor reported no mathematical defect within its
documented scope but retained PUBLICATION PRIVATE HOLD. The author
authorized conditional publication. The auditor's original signed
identity and a verbatim comparison with its chat answer were not
verified. Four optional SAT runs did not prove optimality; they are not
premises of the theorem. No human peer review or proof-assistant
formalization was performed.

\subsection{Prior work and
attribution}\label{prior-work-and-attribution}

Klavžar and Tan (2023, DOI 10.1007/s40840-023-01517-y) study edge
general position in the ordinary Fibonacci and Lucas cubes (k=2); their
coordinate-class construction is prior work. Ilić, Klavžar, and Rho
(2012, DOI 10.2298/AADM120108002I) define generalized Lucas cubes via
avoidance of a cyclic forbidden word: the present graph H(n,k) is their
generalized Lucas-cube model with forbidden word 1\^{}k. The model and
its naming are therefore prior work. Mollard (2020, arXiv:2004.10198)
concerns perfect codes in Lucas and generalized Lucas cubes, a different
graph invariant. The literature review is partial; priority of the
edge-general-position formula and classification is not established.

\begin{center}\rule{0.5\linewidth}{0.5pt}\end{center}

\section{Frozen supplement v0.2: complete source
text}\label{frozen-supplement-v0.2-complete-source-text}

\section{Exact edge general position in cyclic run-avoiding cubes for
every k \textgreater=
3}\label{exact-edge-general-position-in-cyclic-run-avoiding-cubes-for-every-k-3-1}

No.78 v0.2 --- PRIVATE HOLD --- 2026-11-19 JST Author-side
computer-assisted proof candidate. Independent AI audit and Windows
reproduction pending. NOVELTY INCONCLUSIVE / REFERENCES PARTIAL. Not
peer reviewed. All rights reserved.

\subsection{Scope and relationship to frozen
v0.1}\label{scope-and-relationship-to-frozen-v0.1}

This supplement extends the theorem of the byte-frozen v0.1 manuscript
(in FROZEN\_V01.zip) from k\textgreater=4 to every integer
n\textgreater=k\textgreater=3. The parent contains the independent
elementary star and projection lemmas, the count and density
recurrences, and the analytic k\textgreater=4 proof. No result from
No.75--77 is assumed. The new k=3 upper bound uses an exact finite
transfer-matrix certificate and an analytic induction in the gaps, not
enumeration of all large graphs.

Let H(n,k) be the induced graph of the n-cube on binary words with no
cyclic block of k ones; let C\_n\^{}(k) denote its order and A\_m\^{}(k)
the linear run-avoiding word count, with A\_0=1. Edge general position
means that no geodesic contains three distinct selected edges. Set
B\_n=2(C\_n\textsuperscript{(k)-A\_(n-1)}(k)).

\textbf{Theorem (candidate).} For every n\textgreater=k\textgreater=3,
gp\_e(H(n,k))=B\_n, and its maximum sets are exactly the unions of two
distinct complete coordinate edge classes.

The lower bound and equality classification within sets using at most
two directions follow from v0.1 §1. The full case k\textgreater=4 is
v0.1 §§2--5, unchanged. It remains to prove k=3, including the equality
case n=4.

\subsection{1. Four-coordinate exclusion lemma for cyclic triple
avoidance}\label{four-coordinate-exclusion-lemma-for-cyclic-triple-avoidance}

Let n\textgreater=4, let four distinct marked coordinates be given in
cyclic order, and let c be ANY four-bit pattern (no feasibility
restriction is imposed). Let E denote the number of legal cyclic words
differing from c at exactly two of the marked coordinates, and let C
denote the total number of legal words. Then

\begin{verbatim}
11 E >= 3 C.
\end{verbatim}

Proof. Let the positive integer gaps g\_1,\ldots,g\_4 be the cyclic
distances between successive marked coordinates, summing to n.~A
rotation lets the first marked coordinate be 0. The state r=0,1,2 counts
the trailing ones. With rows denoting source states and columns denoting
destination states, write

\begin{verbatim}
T_0 = [[1,0,0],[1,0,0],[1,0,0]],
T_1 = [[0,1,0],[0,0,1],[0,0,0]],
T = T_0+T_1.
\end{verbatim}

For a prescribed marked bit pattern b, the number of legal completions
is

\begin{verbatim}
N_g(b) = tr( T_(b_1) T^(g_1-1) ... T_(b_4) T^(g_4-1) ).
\end{verbatim}

Indeed each closed state walk gives a cyclic word with no triple ones.
Every legal word has at least one zero, which forces state 0 immediately
after it and determines all states uniquely; conversely the all-one word
cannot give a closed walk. This proves the trace identity, with fixed
labeled positions and no division by rotations. Thus C=sum\_b N\_g(b)
and E=sum\_(d\_H(b,c)=2) N\_g(b).

Direct matrix multiplication gives T\textsuperscript{3=T}2+T+I. For any
g\_i\textgreater=4, the value L(g,c)=11E-3C is the sum of its three
values with g\_i replaced by g\_i-1,g\_i-2,g\_i-3 and all other gaps and
c unchanged. This follows by linearity in the corresponding power of T,
without any positivity assumption on the coefficients in L. All new gaps
remain positive, so lowering a gap also leaves a valid four-mark cyclic
instance. Induction on the total gap sum reduces L\textgreater=0 to g in
\{1,2,3\}\^{}4 and all c in \{0,1\}\^{}4. This inclusion of ALL c is
essential: no assumption about preserving feasibility under gap
reduction is made.

The finite certificate verify\_v02.py verifies 81 gap tuples times 16
center patterns = 1,296 inequalities by exact integer matrix
multiplication. It independently counts each projection histogram by
directly enumerating all words of the corresponding length (4 through
12). Both methods agree on every histogram. Minimum E/C is 3/11 and
there are eight base equality cases. This completes the
computer-assisted induction. The lower bound 3/11 is sharp as a
universal four-coordinate inequality (for example n=4, center pattern of
weight one).

\subsection{2. Upper bound for k=3,
n\textgreater=5}\label{upper-bound-for-k3-n5}

Suppose S uses at least three directions. By the elementary star lemma
proved in v0.1 §1, it is a star forest. If its largest star degree
d\textless=3, then \textbar S\textbar{} \textless= 3C/4 \textless{}
19C/25 \textless= B\_n, using the density recurrence proved in v0.1 §2.

If d\textgreater=4, choose any four arms of a star at vertex z. Let c be
its projected pattern on these four coordinates. Every legal vertex x
counted by E in §1 has two agreements and two disagreements with z on
these coordinates. For any incident selected edge direction l, remove l
from the four marked coordinates if it is present. At least one
agreement i and one disagreement j remain. These are two different arms
of the star, distinct from l. The projection of the incident edge onto
i,j is a leaf of the two-arm star. The projection lemma of v0.1 §1
forbids that edge. Thus all E vertices are excluded from incidence with
S. Since a star forest has no more edges than its incident vertices,
\textbar S\textbar{} \textless= C-E \textless= 8C/11 \textless{} 19C/25
\textless= B\_n.~Both degree cases are strict. This proves the value and
classification for all n\textgreater=5, k=3.

\subsection{3. Small cases and equality}\label{small-cases-and-equality}

For n=3, k=3: C=7, B=6, and any S using at least three directions has
\textbar S\textbar\textless=3C/4\textless6. For n=4, k=3: C=11, B=8. The
degree bound gives \textbar S\textbar\textless= floor(4C/5)=8,
irrespective of star degree for sets using at least three directions. To
settle uniqueness of the coordinate-pair form even here, verify\_v02.py
separately builds the actual 11-vertex, 16-edge graph, computes BFS
distances from every vertex, and enumerates every geodesic by moving
along edges of increasing source distance. It generates 96 forbidden
edge triples, tests all binomial(16,8)=12,870 eight-edge subsets and all
binomial(16,9)=11,440 nine-edge subsets. Exactly six eight-edge sets
survive, precisely the unions of two complete coordinate classes; no
nine-edge set survives. Any larger set would have a nine-edge subset, so
this also independently confirms the value 8.

This completes the stated candidate theorem within
n\textgreater=k\textgreater=3. No n\textless k convention is asserted.
No k=2 result is assumed or claimed by this supplement.

\subsection{Verification and
limitations}\label{verification-and-limitations}

Parent v0.1 manifest, finite verifier, and expected stdout checked
afresh from the frozen ZIP. The new exact verifier uses Python 3.10+
standard library, no numerical solvers. Its finite checks are part of
the k=3 proof and do not independently establish the analytic lemmas.
Independent mathematical audit of the complete parent and supplement is
still required. The source search is inherited as preliminary only; no
additional priority claim is made in this revision. See parent
references. Do not publish.

\begin{center}\rule{0.5\linewidth}{0.5pt}\end{center}

\section{Frozen parent v0.1: complete source
text}\label{frozen-parent-v0.1-complete-source-text}

The PRIVATE HOLD and unresolved k=3 statements below are historical
statements from the frozen parent. Consult the v0.2 supplement above for
the full final mathematical scope and the author's later release
authorization.

\section{Exact edge general position numbers of cyclic run-avoiding
cubes for k \textgreater=
4}\label{exact-edge-general-position-numbers-of-cyclic-run-avoiding-cubes-for-k-4}

No.78 v0.1 --- PRIVATE HOLD --- 2026-11-19 JST Author-side analytic
proof candidate. Independent AI audit and Windows reproduction pending.
NOVELTY INCONCLUSIVE / REFERENCES PARTIAL. Not peer reviewed. All rights
reserved.

\subsection{Model and results}\label{model-and-results}

For integers n\textgreater=k\textgreater=3, let H(n,k) be the induced
subgraph of the binary n-cube on words containing no cyclic block of k
ones. Write C\_n=C\_n\^{}(k) for its order, and A\_m=A\_m\^{}(k) for the
number of linear length-m words avoiding k consecutive ones (A\_0=1). A
selected edge set S is in edge general position if no three distinct
selected edges lie on a common geodesic. Let B\_n=2(C\_n-A\_(n-1)).

Theorem: For all n\textgreater=k\textgreater=4, gp\_e(H(n,k))=B\_n.~The
maximum sets are exactly the unions of any two distinct complete
coordinate edge classes.

Partial result for k=3: gp\_e=B\_n at n=3,4. For n\textgreater=5, any S
larger than B\_n must use at least three coordinates and have largest
star degree between 4 and 11 inclusive. No claim of equality for all k=3
is made.

No.75--77 main theorems are not assumed. The elementary projection and
star lemmas used in No.77 are reproduced below. The proof of the theorem
is analytic; finite checks are corroboration only.

\subsection{1. Geometry and lower bound}\label{geometry-and-lower-bound}

The legal family is downward closed. A path via x AND y between any
legal x,y attains Hamming distance, so H(n,k) is isometric in Q\_n.~A
geodesic never repeats a coordinate. Each coordinate class has
C\_n-A\_(n-1) edges: a legal word with a one there corresponds to its
downward edge, and words with a zero there are exactly linear legal
words of length n-1. Rotation gives the same count in every direction.
Any union of two classes is in edge general position, of size B\_n.~A
set using at most two coordinates cannot exceed B\_n; equality requires
two complete classes.

Projection lemma: If S contains two edges of a star at c, in coordinates
i,j, then a selected edge in another coordinate l cannot project on i,j
to a leaf of that two-arm star. Indeed let z be the corresponding leaf
and w the endpoint of the edge with w\_l=z\_l. These vertices agree on
i,j,l. Concatenate the other leaf -\textgreater{} c -\textgreater{} z, a
geodesic from z to w, and the selected edge. The directions of these
parts are disjoint, so this is a geodesic containing three selected
edges, a contradiction.

Star lemma: If S uses at least three directions, the graph of its
selected edges is a star forest. A simple three-edge path must repeat a
direction, else it is a geodesic. Its directions must then be i,j,i. The
two consecutive two-arm stars have leaf projections on i,j covering all
four patterns, so the projection lemma forbids a third direction
anywhere. Thus no simple three-edge path exists. The selected graph is
bipartite, hence its nontrivial components are stars. If the maximum
degree is d and E vertices are excluded from incidence with S, then
\textbar S\textbar\textless=d/(d+1)(C\_n-E).

\subsection{2. Counting and density}\label{counting-and-density}

The linear counts satisfy A\_m=2\^{}m for 0\textless=m\textless k and
A\_m=sum\_\{j=1\}\^{}k A\_(m-j) for m\textgreater=k. The cyclic counts
satisfy C\_j=2\^{}j-1 for 1\textless=j\textless=k under the
periodic-word convention at j\textless k, and C\_n=sum\_\{j=1\}\^{}k
C\_(n-j) for n\textgreater k. To see this, take the k-state matrix T of
transitions of the run of ones: every state r=0..k-1 can move to state 0
by emitting zero, and r\textless k-1 can move to r+1 by emitting one.
Legal cyclic words correspond bijectively to closed walks with a fixed
starting index, since at least one zero determines all run states
uniquely. Thus C\_j=tr(T\^{}j) for positive j. The characteristic
polynomial is x\textsuperscript{k-x}(k-1)-\ldots-1 (companion matrix, up
to transposition/convention), and Cayley-Hamilton gives the recurrence
for n\textgreater k. Consequently any fixed linear combination of C\_n
and A\_(n-1) satisfies this recurrence for n\textgreater=k+1, and a
nonnegative (or positive) block of k consecutive values propagates
forward.

For k\textgreater=5, B\_n/C\_n\textgreater10/11 for n\textgreater=k.
This is equivalent to 6C\_n-11A\_(n-1)\textgreater0. At n=k this
difference is 2\^{}(k-1)-6\textgreater0. For n=k+t with
1\textless=t\textless=k-1, two disjoint starts of maximal runs of length
at least k in a non-all-one cyclic word are impossible (would require at
least 2k+2 positions). A bad non-all-one cyclic word has a unique
occurrence of 0 followed by k ones. Hence
C\_n=2\textsuperscript{n-1-n\emph{2\^{}(n-k-1). For linear length m=n-1,
an illegal word similarly has exactly one maximal long-run start, either
at the first position or preceded by zero. Thus
A\_(n-1)=2\^{}(n-1)-(t+1)}2}(t-2). Substitution gives
6C\_n-11A\_(n-1)=2\textsuperscript{(t-2)(2}(k+1)-12k-t+11)-6. For t=1
the expression in parentheses is at least 14 (at k=5, increasing in k),
hence the difference is at least 1. For 2\textless=t\textless=k-1 it is
at least 2\^{}(k+1)-13k+12\textgreater=11 (increasing for
k\textgreater=5), hence the difference is \textgreater=5. The k positive
bases n=k..2k-1 propagate by recurrence.

For k=4, B\_n/C\_n\textgreater=11/13 for n\textgreater=4: verify
n=4,5,6,7 by C=(15,26,51,99), A\_(n-1)=(8,15,29,56), and propagate
13B\_n-11C\_n.~In addition B\_n/C\_n\textgreater6/7 for n\textgreater=6:
verify n=6,7,8,9, yielding C=(51,99,191,367), A\_(n-1)=(29,56,108,208),
hence 7B\_n-6C\_n=(2,8,16,24), all positive. (At n=9:
7\emph{318-6}367=24.) Propagate this block of four values.

For k=3, B\_n/C\_n\textgreater=19/25 for n\textgreater=5: the three
bases n=5,6,7 have C=(21,39,71), A\_(n-1)=(13,24,44), B=(16,30,54). The
differences 25B-19C are (1,9,1), and the recurrence propagates them.

\subsection{3. Uniform vertex exclusion}\label{uniform-vertex-exclusion}

Suppose a star at c has d arms. Choose m pairwise nonadjacent cyclic
coordinates among the arms. The induced graph of these coordinates in
the cycle is a union of paths if proper, or a cycle if full. Thus at
least floor(d/2) independent coordinates exist; for d odd and the
coordinate set proper this improves to ceil(d/2).

Let X be a uniform legal vertex. Conditional on specified bits at any
subset of the chosen coordinates, the probability of either value at any
other chosen coordinate is between 1/5 and 4/5. Deleting a one injects
the conditioned one class into the zero class. Turning a zero to one and
clearing both cyclic neighbors gives a legal word (the new one is
isolated; downward closure handles the other positions) and preserves
the condition. It has at most four preimages, since only the two former
neighbor bits are unknown. Thus the zero class is at most four times the
one class. The same bounds apply to agreement or disagreement with the
fixed center c.

Let Y be the number of disagreements on these m coordinates. The event
2\textless=Y\textless=m-2 excludes the vertex from every selected edge:
after ignoring its incident edge direction, at least one agreeing and
one disagreeing star coordinate remain. Their two-bit projection is a
leaf of a two-arm star, forbidden by the projection lemma.

A uniform lower bound p\_m on this event holds even for adaptive
conditional distributions. Set f\_m(s)=1 for 2\textless=s\textless=m-2
and 0 otherwise. For j=m-1,\ldots,0 set f\_j(s)=min\{(4
f\_(j+1)(s)+f\_(j+1)(s+1))/5,(f\_(j+1)(s)+4 f\_(j+1)(s+1))/5\}. Backward
induction using the conditional probability interval proves the bound
p\_m=f\_0(0). Exact arithmetic gives p\_4=9/125, p\_5=97/625,
p\_6=757/3125. In particular \textbar S\textbar{} \textless=
d/(d+1)(1-p\_m) C\_n, and also \textbar S\textbar{} \textless= (1-p\_m)
C\_n.

\subsection{4. Upper bound for
k\textgreater=5}\label{upper-bound-for-k5}

Assume three or more selected directions. Then S is a star forest. If
d\textless=10, \textbar S\textbar\textless=10 C\_n/11\textless B\_n.~If
d\textgreater=11, at least five independent coordinates exist, so
\textbar S\textbar\textless=(528/625)C\_n \textless10
C\_n/11\textless B\_n.~This gives strict inequality for all sets using
three or more directions.

\subsection{5. Upper bound for k=4}\label{upper-bound-for-k4}

Again assume three or more directions. When d\textless=5,
\textbar S\textbar\textless=5 C\_n/6\textless11
C\_n/13\textless=B\_n.~If d=6 then n\textgreater=6, where
B\_n/C\_n\textgreater6/7, giving strict inequality. If
7\textless=d\textless=9 and n\textgreater=8, there exist four pairwise
nonadjacent arm coordinates (in particular for d=7 it is a proper subset
of the cycle). Hence \textbar S\textbar\textless= (9/10)(116/125) C\_n =
(522/625) C\_n \textless{} (11/13) C\_n \textless= B\_n.~If
d\textgreater=10, use five independent arms and
\textbar S\textbar\textless= (528/625) C\_n \textless{} (11/13) C\_n
\textless=B\_n.

The sole missing case is n=d=7. All seven directions are arms of a star.
Its center c is legal and has Hamming weight at most five, since six
ones on a 7-cycle contain a run of six. If weight(c)\textgreater=2, the
all-zero word has at least two agreements and two disagreements with
c.~If weight(c)\textless=1, select two zero positions of c and set just
those two positions to one, and all other positions to zero; this word
is legal for k=4 and has between two and three disagreements and at
least four agreements. The projection lemma excludes the resulting
vertex. Thus \textbar S\textbar\textless=7*98/8=85.75 \textless{}
B\_7=86. This handles all cases, including equality, and proves the
theorem.

\subsection{6. Partial result for k=3}\label{partial-result-for-k3}

At n=3 all sets using three directions satisfy
\textbar S\textbar\textless=3\emph{7/4\textless6=B\_3. At n=4 the degree
bound yields \textbar S\textbar\textless=floor(4}11/5)=8=B\_4. At
n\textgreater=5, d\textless=3 gives \textbar S\textbar\textless=3/4 C\_n
\textless19/25 C\_n\textless=B\_n.~For d\textgreater=12 choose six
independent arms, yielding \textbar S\textbar\textless=2368/3125
C\_n\textless19/25 C\_n\textless=B\_n.~The remaining d=4..11 are not
resolved by these inequalities. No uniqueness claim at n=4.

\subsection{References and scope}\label{references-and-scope}

Klavzar and Tan, Edge General Position Sets in Fibonacci and Lucas
Cubes, Bull. Malays. Math. Sci. Soc. 46 (2023), DOI
10.1007/s40840-023-01517-y. Ordinary k=2 problem and standard two-class
construction. That construction is not claimed new. Mollard, The
(non-)existence of perfect codes in Lucas cubes, arXiv:2004.10198:
generalized Lucas cube context, different parameter (perfect codes). The
literature search was preliminary; other uses of k-Lucas cube or
generalized Lucas cube can describe different models. Full citation
review and priority assessment pending. No world-first claim. No public
release authorized.

\end{document}
